%%% Related work --- 
%% 
%% Filename: related_work.tex
%% Author: Jacob Ringfjord & Max Wilén
%%% Code:

\chapter{Related work}
\label{cha:related_work}

With the emergence of more advanced tools, binary analysis has witnessed a surge 
in research activity. Given the primary focus of this thesis on 
categorizing the software changes between different application versions, the related work section highlights 
published work with this shared focus. To provide clarity and 
structure, the section is organized based on the work's specific area.


\section{Automated Change Detection and Categorization of Software Versions}
\label{sec:related_work:kandidat_pipeline}
In a previous bachelor's thesis \cite{nfc_kandidatrapport_jesper}, an automated pipeline comparing differences between Android software versions has been implemented. The pipeline uses the Ghidra headless API for binary analysis, and binary diffing is achieved by executing BinDiff on the analyzed binaries. The last step is a classification performed on function names, where the differences are categorized into important or unimportant changes. A visual representation can be seen in Figure \ref{related_work:workflow}.

In this thesis, the classification is the main focus. The previous study \cite{nfc_kandidatrapport_jesper} made an automated pipeline that provides the differences between software versions, but the classification can be significantly improved. Instead of only looking at the names of the outer functions and comparing them with a small set of keywords to determine their importance, this thesis will take a more in-depth approach regarding the classification and provide more details to make it more useful in digital forensics.


\begin{figure}[ht]
    \centering
    \begin{dot2tex}[dot, options=-tmath --autosize, scale=0.9]
    digraph G {
        rankdir=LR;
        node [shape=box, style=filled, fillcolor=lightgray];

        "\text{Old APK}" -> "\text{Binary analysis}" -> "\text{Binary diffing}" -> "\text{Categorization}" -> "\text{Result}";
        "\text{New APK}" -> "\text{Binary analysis }" -> "\text{Binary diffing}";
    }
    \end{dot2tex}
    \caption{Diagram of the automated pipeline for detecting and categorizing software version changes.}
    \label{related_work:workflow}
\end{figure}



\section{Automating binary analysis}
\label{sec:related_work:automating_binary}

In a survey performed by \textcite{automating_binary_analysis}, the realm of binary analysis automation has been categorized into two major categories, namely Data-driven binary analysis and Software-engineering-based binary analysis. Because of this thesis topic, only Software-engineering-based binary analysis will be of interest here. \textcite{automating_binary_analysis} further splits up Software-engineering-based binary analysis into three sub-categories:
\begin{enumerate}
    \item \textit{Offline analysis}: Performing analysis without executing the program
    \item \textit{Online analysis}: Performing analysis by executing the program with given values
    \item \textit{Hybrid analysis}: A combination of offline and online analysis
\end{enumerate}

Leveraging the insights provided by \textcite{automating_binary_analysis} offers a valuable head start for the experimentation process, due to the comprehensive and detailed list of tools presented. As a result, the following sections discuss tools and techniques mentioned in \cite{automating_binary_analysis}, which also closely align with the objectives of the thesis.

\subsection{Offline Analysis}
This static analysis process creates graph representations, such as control-flow graphs, from binary code. Leveraging this higher-level format, \textcite{automating_binary_analysis} highlights various tools with distinct features that facilitate information extraction. CoP \cite{CoP} is one of these, which uses a binary-oriented, obfuscation-resilient approach that looks at the "\textit{longest common subsequence of semantically equivalent basic blocks}" \cite{CoP}. This enables the detection of similar paths even under obfuscation when looking at two different blocks of binaries. 

Another one is XPMChecker \cite{XMPChecker}, which utilizes backward slicing to track instructions executed from specific inputs. From this, it would be possible to reveal the logic within the program that distinguishes between the two software versions. 

Given both of these tools' clear alignment with the objectives of this thesis, CoP and XMPChecker hold significant potential to advance the experimental part of the research.


\subsection{Online Analysis}
This dynamic analysis process captures run-time information that is later used by specialized analysis tools. One of the tools that \textcite{automating_binary_analysis} mentions in the survey is BinSim \cite{binsim}, which performs analysis on the instructions called by the program. More specifically, BinSim uses symbolic execution to detect semantic differences in how instructions (based on opcodes) are utilized when comparing two versions of the same program. As a result, BinSim enables the verification of conditional equivalence between the two versions, essentially determining whether two executables have a common logical behavior.


\subsection{Hybrid Analysis}
This sub-category combines offline and online analysis to create a comprehensive binary analysis approach for extracting information about executables. This approach leverages both static analysis and dynamic tracing, subsequently comparing the retrieved results to gain deeper insights into an executable's behavior. Such comparisons may reveal hidden execution paths that are only accessible at runtime. To achieve this, \textcite{automating_binary_analysis} suggests angr \cite{angr_original} as a viable option. A tool that can \textit{a)} explore execution paths both with symbolic and concrete values (a combination called Concolic testing \cite{concolic_testing}), \textit{b)} capture runtime data, and \textit{c)} perform trace analysis on the captured data. However, the caveat is performance overhead due to path explosion and problems dealing with self-modified code.



\begin{comment}
\section{Improving binary analysis tools}
\label{sec:related_work:improving_binary_tools}
The reverse engineering tool Ghidra \cite{ghidra_official}, also mentioned in section \ref{sec:background:binary_disassembler_decompilers}, can be of great use when performing static analysis on 
software, enabling symbolic execution as part of the program
would entail increased computational power and improved 
functionality. This is something that 
\textcite{flum_huber_ghidrion_2024} have provided work towards, 
by creating a Ghidra plugin that makes use of symbolic 
execution tools by integrating them into Ghidra. This makes 
the process of analyzing differences between software versions 
more streamlined due to the use of dynamic analysis. For instance, finding 
the different execution paths of a function are no longer 
needed to be done with another tool. To summarize, this 
integration of symbolic execution into Ghidra would enable the 
possibility to build complete Control-flow graphs within 
Ghidra, something that previously had to be done with tools 
such as angr \cite{angr_original}.

\textbf{TODO}: A bindiff and ghidra plugin that exports the data needed from a supported disassembler (like ghidra) to run bindiff on it. 
disassemble in ghidra -> export to bindiff format using BinExport -> Use BinDiff on the exported data

Experimenting with a variety of binary analysis tools is undoubtedly
an effective approach to building a robust pipeline for binary analysis. 
Due to this, the project aims to explore multiple binary analysis tools 
and leverage prior research to develop an effective binary analysis tool.
\end{comment}

\section{Android Class Structure Matching}
\label{sec:related_work:class_matching}
\textcite{android_matching} have implemented a tool named ApkDiff for matching obfuscated classes between different Android application versions. ApkDiff uses a bytecode-based matching of the classes to determine which classes relate to each other between versions. This is achieved by comparing semantic features like fields, methods, and call graphs.

Obfuscated code, such as renaming symbols or class names, can protect an application by replacing meaningful names with shortened, random, or misleading ones \cite{android_matching}. For instance, an actual class name is useful for developers and reverse engineers since it indicates what the class is used for, but the Android Runtime does not rely on the names for execution \cite{android_matching}. 

In this thesis, obfuscated code will not be covered. Instead, the result relies on what the chosen binary analysis and diffing tool manage to extract. This thesis categorizes the actual code differences between two software versions and not only the class names. The research conducted by \textcite{android_matching} shares a similar goal with this thesis in reducing the time and effort required for manual analysis.  

\section{Software Version Change Classification}
\label{sec:related_work:classify_detection}
\textcite{classification_of_change_messages} presented a method to classify changes from version control based on change messages. The developed classifier program identifies whether a change is a bug fix, a new feature, or general maintenance by detecting predefined keywords in the commit messages. 

While understanding the type of change could be interesting for this thesis, it is assumed that only the binary code is available and that version control is not public. Therefore, this aspect falls outside the scope of this thesis. However, the approach of labeling changes based on keywords and then categorizing aligns with the methodology used in this thesis. This thesis can draw inspiration from how the results were gathered and the methods used by \textcite{classification_of_change_messages} to present them.   


\begin{comment}
\section{dump}
If applicable - write about previous work who have classified software changes. For instance, new feature, deprecated feature, bug fix, security path, general maintenance, etc. However, most articles do it by machine learning and looking at change metrics from version control rather than code metrics.

Some work has been done regarding classifying malicious and benign android applications. Usually uses some datasets trained for the targeted area. 
\end{comment}




\begin{center}
    * * *\\
    \Large\textit{Summarized Findings}
\end{center}

\noindent
The literature study revealed no prior work that directly categorizes code changes between two software versions in the way this thesis aims to. The closest related work is the bachelor project by \textcite{nfc_kandidatrapport_jesper}, which introduced an automated pipeline for identifying differences between Android software versions. However, its classification method was limited to function names, and the categorization used only addressed whether the changes were of importance or not. This thesis aims to develop a more comprehensive categorization, making it more valuable for digital forensics and reverse engineering. 

While \textcite{android_matching} compared Android application versions and explored the challenge of obfuscated code, which this thesis does not address. Their work shares a similar goal: reducing the effort required for manual analysis of software changes. Additionally, methodologies from change classification research, such as \textcite{classification_of_change_messages}, provide valuable insights into labeling and categorization techniques, even though they use available version control messages and not starting from the binaries.

By building upon these previous works, this thesis aims to bridge the gap in existing research and contribute with a systematic approach to categorizing software changes by analyzing binaries.
